\lecture{15}
\title{Introduction to Hypothesis Testing}

\begin{document}

\subsection{Hypothesis Testing}

\textbf{Definition.}
A \textbf{test} is a function $\Psi: \mathrm{data} \to \{0, 1\}$.
Equivalently, a test is an estimator of a binary value.
\begin{itemize}
    \item The \textbf{rejection region} of a test $\Psi$
    is the set $\Psi^{-1}(1)$ of all data for which $\Psi(\mathrm{data}) = 1$.
    \item A \textbf{test statistic} is a function that summarizes the data
    and is sufficient for computing a test $\Psi$.
\end{itemize}

In other words, a test is a process of reading in data
and answering a question about the population based on that data.
Given a parametric model $\{\mathbb{P}_{\theta} \mid \theta \in \Theta\}$,
the question a test will answer must look like
``Does $\theta \in \Theta_0$ or $\theta \in \Theta_1$?'',
where $\Theta = \Theta_0 \sqcup \Theta_1$.

\textbf{Definition.} The \textbf{null hypothesis} $H_0$ is the claim that $\theta \in \Theta_0$,
and the \textbf{alternative} $H_1$ is the claim that $\theta \in \Theta_1$.

The distinction is that the null hypothesis is the ``status quo'' (e.g.
``the placebo outperforms the drug'').

\subsection{Error Types}

Of course, a test $\Psi$ only reads in sample data,
so the answer that it gives won't always be correct.

\textbf{Definition.} A \textbf{Type I} error occurs
if $H_0$ is true but $\Psi = 1$, and
a \textbf{Type II} error occurs
if $H_1$ is true but $\Psi = 0$.

Type I is conventionally the more serious error:
it is worse to wrongly declare ``the drug outperforms the placebo''.
Note that the probability of a Type I / Type II error
depends on the true population parameter $\theta$.

\textbf{Definition.} The \textbf{size} of a test $\Psi$ is
$\mathrm{size}(\Psi) := \max_{\theta \in \Theta_0} \mathbb{P}_{\theta}(\Psi = 1)$:
the maximum possible probability of a Type I error.
A test $\Psi$ is said to have \textbf{level} $\alpha$
if $\mathrm{size}(\Psi) \leq \alpha$.

\textbf{Definition.} The \textbf{power} of a test $\Psi$ is the function
$\beta(\theta) = \mathbb{P}_{\theta}(\Psi = 1)$.

\begin{center}
    \frame{\includegraphics[width=0.35\textwidth]{lecture-15/power}}
\end{center}

In the above, the size of $\Psi$ is the $y$-intercept $\alpha$.
A perfect test would have $\beta(\theta)$ be the
\href{https://mathworld.wolfram.com/HeavisideStepFunction.html}{Heaviside step function}.

\end{document}